\section{おわりに}
本稿では，画像中の物体らしい領域を検出するタスクと，その領域に対応するカテゴリをアノテーションするタスクを明確に分離する物体アノテーション手法を提案した．
物体領域の検出にはカテゴリに依存しない高精度なバウンディングボックスを1回の推論で出力するDETRを用い，カテゴリのアノテーションには検出済みの領域のクロップに対するVLMのゼロショット分類を用いた．
この分離により，一度検出した物体領域に対し，候補語彙集合を変更するだけで異なる粒度・異なる語彙でのカテゴリのアノテーションを行うことができ，WordNetの上位語・下位語関係を用いることで，追加学習を一切行わずに粒度可変のアノテーションを実現した．

一連の定量評価を通じて，以下を明らかにした．
第一に，オープンボキャブラリなCLIPが閉域的なResNet50-IN1kに対して統計的に有意に高い単語推定精度を示し，カテゴリのアノテーションにはオープンボキャブラリなVLMが必要であることを確認した．
第二に，物体領域の検出モデルとしてDETR，OWL-ViT，Grounding DINOを比較し，位置検出精度においてDETRが劣らないことを確認した上で，検出を1回で済ませられるDETRの採用が妥当であることを示した．
第三に，提案パイプライン全体でTop-1正解率0.604，Top-5正解率0.894を達成し，クロップを正方形にパディングしてVLMに入力することでTop-1正解率が約0.08向上することを示した．
また，COCOに含まれない語を主要な被写体とする実画像に対しても追加学習なしに正しくアノテーションできることを確認した一方，学習データと視覚的性質が大きく異なる物体（灯台）に対してはDETRの検出自体が失敗することも確認した．
第四に，WordNet階層を用いた粒度可変アノテーションについて，最初の細分化の段階で確信度が大きく低下すること，および下位語が1語しかない分岐が細分化の約18\%で生じ，深い階層の見かけの確信度を押し上げていることを明らかにした．
第五に，COCOの80クラスで学習したYOLOと同一の語彙で比較し，80クラスの範囲内ではYOLOの精度が大幅に高く，その差は検出器ではなく分類器に起因すること，一方でYOLOはCOCO外の物体を80クラス内の語でしかアノテーションできないことを示した．
第六に，深さ0をYOLOの予測とし，深さ1以降をVLMとWordNetで細分化するハイブリッド構成により，深さ0のTop-1正解率0.981を保ちつつ，粒度可変のアノテーションの対象となるインスタンスを\add{，評価インスタンス数が近い条件で比較して提案手法の1.49〜1.65倍}に増やせることを示した．
\add{第七に，LVISの細かい粒度の正解を用いて，深さ0が正しくても深さ1の推定が正しいのは半数未満であること，誤った枝における深い段階の確信度は誤りの検出に役立たないこと，および累積確信度による細分化の停止によって階層的適合率が提案手法で0.425から0.611，ハイブリッド構成で0.594から0.801へ向上することを示した．}

今後の課題として，WordNet階層のフィルタリング，物体領域の検出モデルの大規模な比較，SAM等の汎用的なセグメンテーションモデルとの組み合わせ，候補語彙に正解語が存在しない領域を棄却する仕組みの導入，\add{および深さ1以降の確信度を正誤の手がかりとして利用できるようにする細分化の手法}が挙げられる．
